\section{Results}
\label{sec:results}

\subsection{Direction Quality Varies by Concept Type}
\label{sec:quality}

\Tabref{tab:direction_quality} reports leave-one-out consistency across the three concept types.
Morphological relations achieve the highest consistency ($0.383 \pm 0.117$), followed by encyclopedic ($0.346 \pm 0.092$) and lexicographic ($0.194 \pm 0.081$).
The difference between morphological and lexicographic directions is large and significant (Mann-Whitney $U = 48.0$, $p = 0.006$, Cohen's $d = 1.88$).

\begin{table}[t]
    \centering
    \caption{Direction quality by concept type. \loo is the mean leave-one-out cosine similarity ($\pm$ std).
    Categories with \loo $< 0.1$ are excluded from composition experiments.
    Morphological relations form the strongest directions; seven of ten lexicographic categories fail the quality threshold entirely.}
    \begin{tabular}{@{}lccc@{}}
        \toprule
        \textbf{Concept Type} & \textbf{\loo (mean $\pm$ std)} & \textbf{$N$ categories} & \textbf{Surviving filter} \\
        \midrule
        Morphological  & $\textbf{0.383} \pm 0.117$ & 17/18 & 17 \\
        Encyclopedic   & $0.346 \pm 0.092$          & 5/5   & 5  \\
        Lexicographic  & $0.194 \pm 0.081$          & 3/10  & 3 (7 excluded) \\
        \bottomrule
    \end{tabular}
    \label{tab:direction_quality}
\end{table}

The strongest individual directions are verb tense transformations: infinitive$\to$3pSg (\loo$= 0.63$), infinitive$\to$V\textsc{ing} ($0.47$), and infinitive$\to$V\textsc{ed} ($0.41$).
Among encyclopedic concepts, thing$\to$color achieves $0.49$.
In contrast, synonyms ($0.02$), antonyms ($-0.04$), and hyponyms ($0.01$) are essentially non-linear.
This is our most important finding: {\bf the question ``which directions are compositional?'' is largely answered by ``which concepts have coherent linear representations in the first place.''}

\Tabref{tab:taxonomy} provides a full taxonomy, grouping the 25~surviving categories and 7~excluded categories by quality level.

\begin{table}[t]
    \centering
    \caption{Compositionality taxonomy. Categories are classified by \loo consistency into four tiers.
    Strong and moderate directions compose reliably; weak directions compose with caveats; non-linear categories are excluded from composition analysis.}
    \begin{tabular}{@{}lccccr@{}}
        \toprule
        \textbf{Quality} & \textbf{Criteria} & \textbf{Morph.} & \textbf{Encyc.} & \textbf{Lex.} & \textbf{Total} \\
        \midrule
        Strong      & \loo $> 0.4$    & 6  & 1 & 0 & 7  \\
        Moderate    & \loo $\in [0.2, 0.4]$ & 10 & 4 & 1 & 15 \\
        Weak        & \loo $\in [0.1, 0.2]$ & 1  & 0 & 2 & 3  \\
        Non-linear  & \loo $< 0.1$    & 0  & 0 & 7 & 7  \\
        \bottomrule
    \end{tabular}
    \label{tab:taxonomy}
\end{table}


\subsection{Block-Diagonal Structure in Causal Inner Products}
\label{sec:block_diagonal}

\Tabref{tab:cip} compares the magnitude of causal inner products within and across concept categories.
Within-category pairs show significantly higher $|\text{CIP}|$ ($89.07 \pm 36.84$) than cross-category pairs ($75.68 \pm 31.56$; Mann-Whitney $U$, $p = 3.3 \times 10^{-4}$, Cohen's $d = 0.39$).
This confirms the block-diagonal structure observed by \citet{park2024linear}: related concepts share subspace structure, supporting hypothesis~H1.

\begin{table}[t]
    \centering
    \caption{Causal inner product magnitudes. Within-category concept pairs show significantly higher values than cross-category pairs, confirming block-diagonal subspace structure.}
    \begin{tabular}{@{}lccc@{}}
        \toprule
        \textbf{Metric} & \textbf{Within-category} & \textbf{Cross-category} & \textbf{$p$-value} \\
        \midrule
        $|\text{CIP}|$ (mean $\pm$ std) & $\textbf{89.07} \pm 36.84$ & $75.68 \pm 31.56$ & $3.3 \times 10^{-4}$ \\
        Cohen's $d$ & \multicolumn{3}{c}{0.39 (small-to-medium)} \\
        \bottomrule
    \end{tabular}
    \label{tab:cip}
\end{table}


\subsection{Composition Fidelity and the Orthogonality Paradox}
\label{sec:cfs_results}

\Tabref{tab:cfs} presents \cfs and interference scores for within-category, cross-category, and random-baseline direction pairs.
Two findings stand out.

\begin{table}[t]
    \centering
    \caption{Geometric composition metrics. \cfs measures preservation of both components under addition; interference measures cross-concept contamination.
    Within-category pairs achieve higher \cfs but also higher interference, revealing a tradeoff.
    The random baseline corresponds to uniformly sampled unit vectors.}
    \begin{tabular}{@{}lccr@{}}
        \toprule
        \textbf{Metric} & \textbf{Within-cat.} & \textbf{Cross-cat.} & \textbf{Random} \\
        \midrule
        \cfs (mean $\pm$ std) & $\textbf{0.735} \pm 0.083$ & $0.706 \pm 0.031$ & $0.707 \pm 0.007$ \\
        Interference          & $0.102 \pm 0.070$          & $0.044 \pm 0.022$ & $\textbf{0.016} \pm 0.003$ \\
        \bottomrule
    \end{tabular}
    \label{tab:cfs}
\end{table}

\para{Within-category \cfs exceeds cross-category.}
This contradicts our hypothesis~H3, which predicted that orthogonal (cross-category) pairs would compose better.
The explanation is geometric: when $\cos(\vd_A, \vd_B) > 0$, the sum $\vd_A + \vd_B$ is \emph{closer} to each component than when the directions are orthogonal, inflating \cfs.
The baseline \cfs for perfectly orthogonal directions is exactly $1/\sqrt{2} \approx 0.707$, and within-category pairs exceed this because their directions are partially aligned.

\para{Interference tradeoff.}
Within-category pairs show $2.3\times$ higher interference ($0.102$ vs.\ $0.044$).
The correlation between \cfs and interference is positive ($r = 0.33$, $p = 3.2 \times 10^{-9}$): the same alignment that boosts geometric fidelity also increases cross-concept contamination.
This reveals two distinct senses of ``compositionality'': geometric preservation (high \cfs) and functional independence (low interference), which cannot be simultaneously maximized.

Additional correlations illuminate the structure:
\begin{itemize}[leftmargin=*,itemsep=0pt,topsep=0pt]
    \item $|\cos(\vd_A, \vd_B)|$ vs.\ \cfs: $r = 0.37$, $p = 3.7 \times 10^{-11}$ (more aligned $\to$ higher \cfs).
    \item \loo vs.\ interference: $r = 0.36$, $p = 1.9 \times 10^{-10}$ (stronger directions produce more interference).
    \item \loo vs.\ \cfs: $r = 0.01$, $p = 0.80$ (direction quality does not predict geometric \cfs).
\end{itemize}


\subsection{Functional Steering Composition}
\label{sec:steering}

\Tabref{tab:steering} reports functional steering accuracy for single and composed directions.
Across all pair types, composed directions $\vd_A + \vd_B$ maintain $>$96\% accuracy on both target concepts.
The maximum accuracy drop from single to composed steering is 3.5~percentage points (within-encyclopedic, concept~A: $1.000 \to 0.965$).

\begin{table}[t]
    \centering
    \caption{Functional steering accuracy. Single-direction accuracy and composed-direction accuracy ($\vd_A + \vd_B$) for both concepts are shown. Composition success is the minimum of composed accuracies.
    Steering composition is robust across all pair types, with $>$96\% accuracy everywhere.}
    \resizebox{\textwidth}{!}{%
    \begin{tabular}{@{}lccccc@{}}
        \toprule
        \textbf{Pair type} & \textbf{Single A} & \textbf{Single B} & \textbf{Composed A} & \textbf{Composed B} & \textbf{Comp.\ success} \\
        \midrule
        Within-morphological   & 1.000 & 0.997 & 1.000 & 0.993 & \textbf{0.997} \\
        Within-encyclopedic    & 1.000 & 1.000 & 0.965 & 1.000 & 0.982 \\
        Cross-morph-encyc      & 1.000 & 1.000 & 0.998 & 0.994 & 0.996 \\
        Cross-morph-lex        & 1.000 & 0.976 & 0.996 & 0.976 & 0.986 \\
        \bottomrule
    \end{tabular}
    }
    \label{tab:steering}
\end{table}

This near-ceiling performance contrasts with the nuanced geometric picture from \secref{sec:cfs_results}.
Despite within-category pairs showing 2.3$\times$ more interference geometrically, the interference magnitude ($\sim$0.1) is small enough that functional steering is not meaningfully degraded.
In the unembedding space, all directions that pass the \loo quality filter are effectively compositional for steering purposes, regardless of whether they are aligned or orthogonal.
